% Andrew Duff: The Twin Executive – Reading Summary
% Introduction to the European Union, Week 7
% Compile: pdflatex twin-executive-reading-summary.tex

\documentclass[11pt,a4paper]{article}
\usepackage[utf8]{inputenc}
\usepackage[T1]{fontenc}
\usepackage[margin=2.5cm]{geometry}
\usepackage{booktabs}
\usepackage{enumitem}
\usepackage{parskip}
\usepackage{hyperref}
\usepackage{xcolor}
\definecolor{eublue}{RGB}{0,51,153}

\hypersetup{colorlinks=true, linkcolor=eublue, urlcolor=eublue}

\begin{document}

\title{Andrew Duff: \textit{The Twin Executive}\\Key Themes \& Arguments --- Rhetorical Analysis}
\author{Introduction to the European Union\\Universidad Complutense de Madrid\\Week 7 --- The Twin Executive (European Council \& Commission)}
\date{}
\maketitle
\vspace{-1em}
\hrule
\vspace{1em}

\section*{Rhetorical Framework (Ethos, Logos, Pathos)}

\begin{table}[h]
\centering
\begin{tabular}{@{}lll@{}}
\toprule
\textbf{Appeal} & \textbf{Meaning} & \textbf{Duff's Use} \\
\midrule
\textbf{Ethos} & Credibility, authority, expertise & Former MEP; Convention participant; institutional reform advocate \\
\textbf{Logos} & Logic, evidence, argument & Treaty citations; institutional analysis; reform proposals \\
\textbf{Pathos} & Emotion, urgency, shared concern & Appeals to democratic deficit; frustration with EU confusion \\
\bottomrule
\end{tabular}
\end{table}

\textbf{Key insight:} Duff combines insider knowledge (Convention experience) with normative argumentation for institutional reform. His \textit{ethos} comes from participation in treaty drafting; his \textit{logos} from systematic analysis of power-sharing; his \textit{pathos} from concern about accountability and clarity.

\section*{Bibliographic Information}

\begin{itemize}[nosep,leftmargin=*]
\item \textbf{Author:} Andrew Duff
\item \textbf{Source:} \textit{Constitutional Change in the European Union} (Palgrave Macmillan, Cham, 2022)
\item \textbf{Chapter:} ``The Twin Executive,'' pp. 9--22
\item \textbf{DOI:} \url{https://doi.org/10.1007/978-3-031-10665-1_2}
\item \textbf{First Online:} 6 September 2022 (Open Access)
\end{itemize}

\section*{Summary \& Key Points}

\subsection*{I. The Problem: No Single Effective Executive}

The EU lacks a single effective executive. Power is split between:
\begin{itemize}[nosep,leftmargin=*]
\item \textbf{European Council} (heads of state/government) --- provides ``impetus'' and ``general political directions''
\item \textbf{European Commission} --- promotes ``general interest,'' proposes legislation, enforces treaties
\end{itemize}

Executive power is ``spread in a fairly complicated, bifurcated way'' between the two. There is no classic separation of powers; leadership and democratic accountability are difficult to locate.

\subsection*{II. The Two Presidents}

Lisbon (2009) created a \textbf{permanent President of the European Council} (2.5-year term, renewable once). This was favoured by Giscard d'Estaing and Tony Blair. But the result is \textbf{bicephalous governance}:
\begin{itemize}[nosep,leftmargin=*]
\item If the two presidents say the same thing $\to$ duplication
\item If they differ $\to$ division
\item Example: Ankara 2021 --- Erdogan seated Charles Michel (European Council) next to him; von der Leyen (Commission) relegated to a sofa
\end{itemize}

\subsection*{III. Commission Under Pressure}

\begin{itemize}[nosep,leftmargin=*]
\item \textbf{Guardian of the treaties:} Commission can bring infringement proceedings (Art. 258 TFEU). Since Barroso (2004), enforcement actions fell steeply.
\item \textbf{Trade-off:} Technocratic legal service (apply rules) vs.\ political arm (extract concessions from European Council).
\item \textbf{Risk:} If Commission president becomes too complicit with European Council decisions, the Commission loses autonomy to uphold EU law.
\end{itemize}

\subsection*{IV. Reform Proposals}

\begin{enumerate}[nosep,leftmargin=*]
\item \textbf{Merge the two presidencies} --- Commission president chairs European Council and General Affairs Council
\item \textbf{Phase out rotating Council presidency} --- Replace with merit-based chairs (2.5 years), like EP committees
\item \textbf{Reduce Commission size} --- 18 members instead of 27 (Lisbon provision not yet applied)
\item \textbf{Reinforce enforcement} --- EU Competition Authority at arm's length; Commission Law Officer; Treasury Secretary; Defence portfolio
\item \textbf{Transparency} --- Publish working documents, agendas, minutes; cameras on in Council chamber
\end{enumerate}

\section*{Main Arguments}

\subsection*{I. European Council vs.\ Commission: Overlapping Competences}

\begin{itemize}[nosep,leftmargin=*]
\item European Council: sets strategic direction, takes binding decisions (e.g.\ carbon targets), arbitrates Council stalemates, budgetary affairs, senior appointments, rule of law, enlargement, treaty revision
\item Commission: proposes legislation, enforces law, manages policies --- but ``much beholden to guidelines set by the European Council''
\item Council of ministers ``clings to executive powers like a limpet'' over sanctions, farm prices, fish quotas --- powers that ``logically reside with the Commission''
\end{itemize}

\subsection*{II. The Lisbon Experiment Has Not Worked}

\begin{itemize}[nosep,leftmargin=*]
\item Van Rompuy (2009--2014): managed financial crisis well
\item Tusk (2014--2019): tension with Juncker over Brexit and foreign policy
\item Michel: less authority; 2019 appointment of von der Leyen was ``messy compromise'' after Spitzenkandidat process collapsed
\item Von der Leyen: narrowly elected; illiberal backsliding in two states; Brexit; COVID; war in Ukraine
\end{itemize}

\subsection*{III. Spitzenkandidat: Promise and Failure}

\begin{itemize}[nosep,leftmargin=*]
\item 2014: Juncker (EPP Spitzenkandidat) nominated and elected --- success
\item 2019: Weber unacceptable to European Council; no common accord between Parliament and Council; von der Leyen appointed without Spitzenkandidat legitimacy
\item Without transnational lists, Parliament's Spitzenkandidaten ``look rather thin''
\end{itemize}

\subsection*{IV. General Affairs Council Underused}

\begin{itemize}[nosep,leftmargin=*]
\item Supposed to prepare European Council and follow up decisions
\item Still chaired by rotating presidency (6 months) --- ``disjointed, variable''
\item Should be chaired by European Council or Commission president (or deputy)
\end{itemize}

\section*{Context \& Analysis}

\begin{itemize}[nosep,leftmargin=*]
\item \textbf{Historical:} Commission descends from ECSC High Authority (Jean Monnet). European Council invented 1974 (Giscard); formalised Lisbon 2009.
\item \textbf{Federalist vs.\ intergovernmental:} Federalists want Commission as core of future EU government; intergovernmentalists resist. Lisbon was ``classic European compromise.''
\item \textbf{Democratic deficit:} Particular cause = ``uneven performance and opacity of the Council as legislator.'' Trilogues often ``black box''; conciliation procedure more open.
\item \textbf{Foreign policy:} President of European Council represents EU externally in CFSP ``at his level''; High Representative (double-hatted) runs day-to-day. ``Foreign governments may be forgiven if they are confused about who speaks for the European Union.''
\end{itemize}

\section*{Relevance to Course}

\begin{itemize}[nosep,leftmargin=*]
\item \textbf{6th Session:} The Twin Executive --- core reading for understanding European Council and Commission
\item \textbf{Institutional design:} Illustrates tension between intergovernmental (European Council) and supranational (Commission) logic
\item \textbf{Treaty evolution:} Lisbon reforms, Spitzenkandidat, Convention debates
\item \textbf{Current debates:} Enlargement, rule of law, Conference on the Future of Europe, treaty revision
\end{itemize}

\section*{Discussion Questions}

\begin{enumerate}[nosep,leftmargin=*]
\item Should the two presidencies be merged? What would be gained and lost?
\item Is the Commission's declining enforcement of EU law (since Barroso) a necessary political adaptation or a dangerous erosion of supranational authority?
\item How does the Spitzenkandidat process affect the legitimacy of the Commission president? Why did it fail in 2019?
\item What role should the General Affairs Council play in coordinating the Council of ministers and servicing the European Council?
\item Does the ``twin executive'' make the EU incomprehensible to citizens and international partners? How might clarity be improved?
\end{enumerate}

\end{document}
